% Gesamttest «Textaufgaben modellieren» — Quelle des PDFs (python3 scripts/build-lp-pdf.py modellieren).
% Musterlösung und Raster: bewertungspaket.tex. Alle Werte mit python3 nachgerechnet (scripts/lp/modellieren/zahlen.py, 08.10.2026).
% Kein \lt/\gt — pdfLaTeX kennt sie nicht, < und > tun es.
% Neufassung 08.10.2026 nach der Prüfung (TODO.md, H1, M3, M4): Teile nach Hilfsmittel. Teil A ohne Taschenrechner
% (RLP 2.3: lineare Systeme «auch ohne Hilfsmittel», bis drei Variablen), Teil B mit Taschenrechner (poly-solv, sys-solv).
% Jede Gleichungsart kommt vor: linear (G2 mit einer Unbekannten, G3), quadratisch (G6), lineares System (G1, G2, G4, G7),
% quadratisches System mit Produkt der Unbekannten (G5). Jede Aufgabe kombiniert Geübtes neu, keine wiederholt ein Modell aus
% Clip, Übung oder Kapitelaufgabe: Ansätze vergleichen (G2), Set über Terme mit einer Unbekannten (G3), Wasser als dritte Sorte
% im System (G4), Eindampfen mit Masse und Anteil unbekannt (G5: Übung «Eindampfen» + Kontrollclip m · p), Abhebung statt
% Einzahlung (G6, in der Übung «Zinseszins» geübt), gleich viel Zins statt Gesamtzins (G7).
\documentclass[11pt]{article}
\usepackage{../lp-druck}
\renewcommand{\lpname}{Textaufgaben modellieren}
\renewcommand{\lpteilgebiet}{2.1 · 2.3}
\renewcommand{\lpbereich}{Grundlagenfach}
\renewcommand{\lpfuss}{Gesamttest · 25 Punkte}
\begin{document}
\lpkopf{Gesamttest}{%
  \vspace{4pt}\noindent\begin{tabularx}{\linewidth}{@{}X X@{}}
  Code (kein Name): \hrulefill & Punkte: \hrulefill\ / 25
  \end{tabularx}}

\begin{lpinfo}{Anleitung}
\textbf{Zeit:} rund 35 Minuten. \textbf{Hilfsmittel:} Teil A \textbf{ohne} Taschenrechner, Teil B \textbf{mit} Taschenrechner (TI-30X Pro).
Schreib zu jeder Aufgabe die \textbf{Deklaration} (Bedeutung und Einheit), den \textbf{Ansatz} und den Rechenweg auf — bewertet wird der Weg.
Wo du den Rechner benutzt, notierst du die Grundform und was du eintippst. Prüfe jede Lösung am Text.
Danach: Lösung scannen oder fotografieren (ohne Namen und Standort) und mit dem \emph{Bewertungspaket} bewerten lassen.
\end{lpinfo}

\lpteil{Teil A · ohne Taschenrechner}{Kapitel 1 und 3 · 11 Punkte}

\begin{lpaufgabe}{G1}{Ziffernrätsel}{4 P}
Eine zweistellige Zahl ist um 9 grösser als das Vierfache ihrer Quersumme. Vertauscht man ihre Ziffern, wird die Zahl um 27 grösser.
Deklariere, stelle das System auf, löse es von Hand und mach die Probe am Text.
\lpplatz{16}
\end{lpaufgabe}

\begin{lpaufgabe}{G2}{Ansätze vergleichen}{4 P}
Ein Theater hat 300 Plätze: im Parkett zu 22 CHF, auf dem Balkon zu 30 CHF. Eine Vorstellung ist ausverkauft und bringt 7400 CHF ein.
Drei Lernende stellen den Ansatz auf, alle mit \(x\): Anzahl Parkettplätze; A und B dazu mit \(y\): Anzahl Balkonplätze.
\begin{itemize}[leftmargin=2.2em,itemsep=0pt]
\item[(A)] \(x + y = 7400\) und \(22 \cdot x + 30 \cdot y = 300\)
\item[(B)] \(x + y = 300\) und \(22 \cdot x + 30 \cdot y = 7400\)
\item[(C)] \(22 \cdot x + 30 \cdot (300 - x) = 7400\)
\end{itemize}
\begin{enumerate}[label=(\alph*),leftmargin=*,itemsep=1pt]
\item Welche Ansätze sind richtig? Was ist beim falschen falsch?
\item Löse einen richtigen Ansatz von Hand und gib an, wie viele Parkett- und Balkonplätze es sind.
\end{enumerate}
\lpplatz{9}
\end{lpaufgabe}

\begin{lpaufgabe}{G3}{Menü und Einzelverkauf}{3 P}
Ein Kiosk verkauft Sandwiches zu 6 CHF, Getränke zu 3 CHF und das Menü «Sandwich + Getränk» zu 8 CHF. An einem Mittag gehen
50 Sandwiches und 70 Getränke weg — einzeln oder im Menü — für zusammen 476 CHF. Wie viele Menüs wurden verkauft?
Deklariere drei Unbekannte, stelle die Gleichungen auf und löse sie von Hand.
\lpplatz{9}
\end{lpaufgabe}

\lpteil{Teil B · mit Taschenrechner}{Kapitel 2 und 4 · 14 Punkte}

\begin{lpaufgabe}{G4}{Zwei Lösungen und Wasser}{3 P}
Aus einer Lösung mit 10\,\% Wirkstoff, einer Lösung mit 30\,\% Wirkstoff und 10 kg Wasser sollen 50 kg Lösung mit 14\,\% Wirkstoff
entstehen. Wie viel kg braucht es von jeder der beiden Lösungen? Deklariere, stelle das System auf, notiere, was du in sys-solv eintippst,
und gib die Antwort.
\lpplatz{11}
\end{lpaufgabe}

\begin{lpaufgabe}{G5}{Eindampfen}{4 P}
Eine Zuckerlösung enthält 6 kg Zucker. Lässt man 10 kg Wasser verdunsten, steigt ihr Zuckeranteil um 5 Prozentpunkte.
Wie schwer war die Lösung, und welchen Zuckeranteil hatte sie? Deklariere, stelle das System auf, bring es in die Grundform
für poly-solv, löse es und begründe, welche Lösung wegfällt.
\lpplatz{13}
\end{lpaufgabe}

\begin{lpaufgabe}{G6}{Zinseszins mit Abhebung}{4 P}
Auf einem Konto liegen 6000 CHF. Nach einem Jahr werden 1000 CHF abgehoben. Der Zins wird jährlich gutgeschrieben und mitverzinst;
der Zinssatz bleibt gleich. Nach zwei Jahren stehen 5110.60 CHF auf dem Konto. Bestimme den Zinssatz: Deklaration, Gleichung,
Grundform, Lösung mit poly-solv — und begründe, welche Lösung wegfällt.
\lpplatz{12}
\end{lpaufgabe}

\begin{lpaufgabe}{G7}{Gleich viel Zins}{3 P}
18\,000 CHF werden auf zwei Anlagen aufgeteilt: Anlage A liegt 9 Monate zu 2\,\%, Anlage B ein ganzes Jahr zu 1.2\,\%.
Beide Anlagen bringen gleich viel Zins. Wie viel CHF liegen auf jeder Anlage, und wie viel Zins bringt jede?
Deklariere, stelle das System auf und löse es mit sys-solv.
\lpplatz{11}
\end{lpaufgabe}
\end{document}
